\section{Implementation details: training protocol}
\label{app:impl-train}
\subsection{Hyperparameters and implementation settings}
\label{app:hyperparams}

\begin{table}[t]
\centering
\footnotesize
\setlength{\tabcolsep}{3pt} % 열 간격 줄이기
\renewcommand{\arraystretch}{1.05}
\caption{Hyperparameters used in the experiments. Values are highlighted for easy modification.}
\label{tab:hyperparams}
\begin{tabular}{p{0.18\linewidth} p{0.36\linewidth} p{0.38\linewidth}}
\hline
\textbf{Category} & \textbf{Setting} & \textbf{Value} \\
\hline
System & System / seed / device & \HP{ogtt\_simul}, \HP{42}, \HP{CUDA if available} \\
Data & \#samples / split / batch & \HP{10000}, \HP{0.2}, \HP{512} \\
SDE data & Aug. factor / diffusion scale & \HP{30}, \HP{1.27} \\
Training & Epochs / lr / wd & \HP{10000}, \HP{$10^{-4}$}, \HP{$10^{-4}$} \\
Training & Early stopping (pat., min-$\Delta$) & \HP{200}, \HP{$10^{-6}$} \\
Training & LR schedule & \HP{Plateau: factor 0.9, pat. 20, min lr $10^{-8}$} \\
Training & Grad clip & \HP{$\|g\|\le 1.0$} \\
Features & Derivative / method & \HP{on}, \HP{spline} \\
Loss & Components & \HP{sup. + rec. (w=1,1)} \\
Inference & Iterations & \HP{10} \\
Model & Hidden dims / activation & \HP{[64,64,64]}, \HP{SiLU} \\
Regularize & Spectral norm / power iters & \HP{on}, \HP{5} \\
\hline
\end{tabular}
\end{table}
\noindent\textbf{Note.}
All parameter-learning and convergence analysis are performed in normalized parameter coordinates.
Reported parameter estimates are denormalized to physical units for evaluation.

This appendix summarizes training and inference hyperparameters used in our experiments.

\subsection{Data and batching}
We generate \HPNumSamples\ trajectories for training, with a test split of \HPTestSplit.
Mini-batch training uses batch size \HPBatchSize.

\subsection{Optimization}
We train for up to \HPEpochs\ epochs using AdamW with learning rate \HPLR\ and weight decay \HPWeightDecay.
We apply gradient clipping with maximum norm $1.0$.

Early stopping is enabled (\HPEarlyStop) with patience \HPEarlyPat\ and minimum improvement threshold \HPEarlyDelta.
A ReduceLROnPlateau scheduler is used as described in the codebase.

\subsection{Loss configuration}
Unless otherwise stated, we use the composite objective with components \HPLossCfg.
The recurrent component corresponds to an unrolled training objective (see Section~\ref{sec:loss-design} for discussion, if included).

\subsection{Model architecture and regularization}
Both networks use hidden dimensions \HPHiddenDims\ and activation \HPAct.
Spectral normalization is enabled (\HPSN) and applied to all linear layers.

\subsection{Derivative features and inference iteration}
Derivative features are used (\HPUseDeriv) and computed at inference by the \HPDerivMethod\ method.
The inference iteration uses \HPIter\ fixed steps.